
\subsubsection{2.1 LLM Calibration and Uncertainty Quantification}

Guo et al. (2017) established Expected Calibration Error (ECE) as the canonical metric for measuring the discrepancy between confidence and accuracy, showing that modern neural networks are systematically overconfident and that temperature scaling can partially recover calibration. However, this work measures calibration globally across all inputs without conditioning on input difficulty.

Kadavath et al. (2022) introduced the P(True) logprob elicitation methodology for LLMs: appending a self-evaluation prompt to the model's own output and extracting the logprob ratio as a confidence signal. They demonstrated that P(True) scales with model capability and achieves meaningful calibration on question-answering tasks. Critically, they study P(True) as a function of model scale rather than problem difficulty.

Yin et al. (2023) investigate LLM self-knowledge, finding substantial gaps in models' ability to identify questions they cannot answer. Their work focuses on factual questions rather than code generation. Liu et al. (2025) and Shorinwa et al. (2024) identify code verification as an open challenge for LLM calibration methods.

\subsubsection{2.2 Code Generation Evaluation}

Chen et al. (2021) introduced HumanEval and the pass@k metric. Austin et al. (2021) contributed MBPP. Liu et al. (2023) released EvalPlus, augmenting HumanEval and MBPP with substantially more test cases (HumanEval+: 80x more tests; MBPP+: 35x more tests), reducing pass@k inflation. None of these benchmark papers analyze calibration quality.

Roziere et al. (2023) introduced Code Llama, a family of code-adapted models fine-tuned from Llama base weights. This model's training data includes a large proportion of common Python utility patterns, a property relevant to the calibration inversion observed in this study.

\subsubsection{2.3 Difficulty Estimation and Stratified Analysis}

Several benchmarks study problem difficulty using external labels (Srivastava et al., 2022; Liang et al., 2022). The self-contained difficulty approach used here, stratifying by each model's own k=5 pass@1 distribution, differs in that it provides a model-relative definition of difficulty. Hard problems are those where a specific model generates 0/5 correct solutions; easy problems are those where it generates 3-5/5 correct solutions.
